\section{Day 31: Big Theorem for Modules (Jan.\ 28, 2026)}
Recall the big theorem\footnote{i know my notes are messy. i'll release a compilation or something before final, dw}: if $M$ is a finitely generated module over a PID $R$, then the following decomposition is unique,
\[ M \cong R^k \oplus \bigoplus R\,/\!\left<p_i^{s_i}\right>, \]
where $k \geq 0$, $s_i \geq 1$, and $p_i$ are prime. Existence has been proven when $R$ is a Eucliean domain. We go over some examples; consider \begin{parlist} \item $V$ a vector space over a field $F$, \item abelian groups over $\ZZ$, \item $R/I$ over a PID $R$, \item $V$ over $F[x]$ if and only if $(V, T \colon V \to V)$ by $Tu = xu$. \end{parlist}
\begin{enumerate}[(i)]
    \item $V$ is isomorphic to $F^k$, so there are no primes.
    \item $A/\ZZ$ is the same as above.
    \item $\Tor(M)$ is given by the collection of elements $m \in M$ where there exists some nonzero $r \in R$ such that $rm = 0$. Then $\Tor(A \oplus B) = \Tor(A) \oplus \Tor(B)$, whence $A \oplus B = \{(a, b) \mid a \in A, b \in B\}$.
    \begin{claim}
        $(a, b)$ is torsion if and only if both $a$ and $b$ are torsion.
    \end{claim}
    \begin{proof}
        In the forwards direction, if there exists nonzero $r \in R$, then $r(a, b) = (0, 0)$ implies $ra = 0$ and $rb = 0$. Conversely, if $r_1 a = 0$ and $r_2 b = 0$, then $r_1 r_2 (a, b) = (0, 0)$.
    \end{proof}
    Thus, $\Tor(R/I) = R/I$ as $I = \left<a\right>$ for some $a \in R$. Similarly, $\Tor(R^k / \left<p_i^{s_i}\right>) = \bigoplus R\,/\!\left<p_i^{s_i}\right>$, so we may write
    \[ R/I = R^k \oplus \bigoplus R\,/\!\left<p_i^{s_i}\right>, \]
    which we apply to $R\,/\!\left<a\right> = R/I \cong \bigoplus R\,/\!\left<p_i^{s_i}\right>$. To see this, write $a = \prod p_i^{s_i}$.
    \item Assume further that $F$ is finite dimensional and algebraically closed. If $f$ is a polynomial in $F[x]$, with $f$ nonconstant, then there exists $\lambda \in F$ such that $f(\lambda) = 0$ (for example, the fundamental theorem of algebra says that the complex numbers are algebraically closed). Suppose $p(x) \in F[x]$ is prime, hence irreducible; if $\det p = 0$, then $p$ is a unit, and is not prime. Thus, $\deg p \geq 1$. This means there exists some $\lambda$ such that $p(\lambda) = 0$, so $(x - \lambda) \mid p$. Indeed,
    \[ p = (x - \lambda) q + r, \]
    where $\deg r < \deg q$ (by the Euclidean algorithm). Substitute $x = \lambda$; then $0 = p(\lambda) = 0 \cdot q(\lambda) + r(\lambda)$ so $r = 0$; thus, up to a unit, $p = x - \lambda$.\footnote{what the helly?} By the big theorem,
    \[ V \cong F[x]^k \oplus \bigoplus F[x]\,/\!\left<(x - \lambda_i)^{s_i}\right>. \]
    Since this is finite dimensional over $F$, we have $k = 0$, so the above consists of only torsion terms, and we get that $\dim V = \sum s_i$.
\end{enumerate}
As an aside, what is $F[x]\,/\!\left<(x - \lambda)^s\right>$ as a module of $F[x]$, i.e., a vector space with a linear transformation $T$?
\begin{claim}
    $\dim_F F[x]\,/\!\left<(x - \lambda)^s\right> = s$, with basis $x^0, \dots, x^{s-1}$.
\end{claim}
Indeed, every $f \in F[x]$ can be written in the form
\[ f = q(x - \lambda)^s + \sum_{i=0}^{s-1} a_i x^i \mod \left<(x - \lambda)^s\right>, \]
so $f = \sum_{i=0}^{s-1} a_i x^i$, i.e., $x^0, \dots, x^{s-1}$ are generators. They are linearly independent, since
\[ 0 = \sum_{i=0}^{s-1} a_i x^i \]
in the quotient implies $(x - \lambda)^s \mid \sum a_i x^i \implies a_i = 0$. Alternatively, we can construct new basis $u_0 = (x - \lambda)^0, \dots, u_i = (x - \lambda)^i$; we have that
\[ (x - \lambda) u_i = \begin{cases} u_{i+1} & i < s - 1, \\ 0 & i = s - 1. \end{cases} \]
Let $T u_i = u_{i+1} + \lambda u_i$ for $i < s-1$ and $T u_{s-1} = \lambda u_{s-1}$; then the matrix of $T$ relative to the basis $\{u_i\}$ is
\[ \begin{pmatrix} \lambda \\ 1 & \lambda \\ & 1 & \lambda \\ & & \ddots & \ddots \end{pmatrix} \]
This is the Jordan canonical form. $A \colon R^\times \to R^s$ leaves $M_A$ unchanged if you do invertible row and column operations on $A$, of all $A$s equivalent to ours, and find one with $a_{11}$ of lowest norm. A general PID is a UFD; find $A$ with $a_{11}$ having a least number of divisors. In a PID, $\gcd(a_{11} = a, b) = g = sa + b$; then
\[ \begin{pmatrix} a & b \end{pmatrix} \begin{pmatrix} s & -b/g \\ t & a/g \end{pmatrix} = \begin{pmatrix} g & 0 \end{pmatrix}. \]
As seen previously, we also have that
\[ \begin{pmatrix} s & -b/g \\ t & a/g \end{pmatrix}\inv = \begin{pmatrix} a/g & b/g \\ -t & s \end{pmatrix}. \]
%We give a short review of tensor products in preparation for next lecture. Recall that $M \oplus N = \{(m, n) : m \in M, n \in N\}$ and $M \otimes N$ is the module generated by symbols of the form $m \otimes n$ (where $m \in M$, $n \in N$), 